%%%%%%%%%%%%%%%%%%%%%%%%%%%%%%%%%%%%%%%%%%%%%%%%%%%%%%%%%%%%
%% CHAPTER 15 -- A Closing Note, and Where to Go Next
%%%%%%%%%%%%%%%%%%%%%%%%%%%%%%%%%%%%%%%%%%%%%%%%%%%%%%%%%%%%

\chapter*{A Closing Note, and Where to Go Next}
\addcontentsline{toc}{chapter}{A Closing Note, and Where to Go Next}

\section*{What you now have}
\addcontentsline{toc}{section}{What you now have}

The preface opened with a promise: by the time you finish,
you will not find these systems mysterious. You may find them
more remarkable than you did when you began, and you will
certainly find them less trustworthy in some specific and
useful ways.

That promise is, we hope, substantially kept.

You now know that a Large Language Model is not a database
that looks up the right answer. It is a function, learned
from text, that produces the most plausible continuation
of a given prefix. Plausibility is its strength. The gap
between plausibility and truth is its characteristic failure.

You know what is inside the model: a deep stack of simple
arithmetic operations, each layer combining the work of the
layers below it, the whole trained by sending errors
backward through the network and nudging each of several
billion weights in the direction that would have reduced
the error. You know why this is expensive, why it required
the coincidence of three things arriving together in the
early 2010s --- the algorithm, the data, and the GPU ---
and why the resulting function, for all its opacity, is
not a black box in principle but a very large grey one.

You know what attention is: the mechanism that lets each
position in a sequence consult every other position,
weighted by a learned measure of relevance. You know that
the weights are not looked up; they are produced by
a learned matching between queries and keys, and that this
matching is what the model discovered, through training,
to be useful for predicting the next token. You know that
the geometry of the embedding space --- the vector space
into which words are placed before the network processes
them --- is not a designed feature but a residue of the
pressure the training imposed: words that share prediction
contexts end up near each other because the cheapest path
to a lower loss is a shared vector.

You know the alignment pipeline: supervised fine-tuning to
establish the format, preference modelling to teach the
model a proxy for human preferences, policy optimisation to
shift the model's output distribution toward higher-scoring
responses. You know that the reward model is a
pattern-matcher rather than a judge, and that this is the
root of sycophancy and reward hacking. You know that
aligned to someone's values is not the same as aligned to
everyone's.

You know the five axes of the error model: structure,
meaning, context, groundedness, environment. You have walked
the ladder on a specific example. You have, in the previous
chapter, a set of questions that put this vocabulary to
work.

You know the lineage: Markov in 1913, Shannon in 1948,
McCulloch and Pitts in 1943, Rosenblatt in 1958, the first
AI winter, the thaw with backpropagation, Hochreiter and
Schmidhuber's long short-term memory in 1997, Bengio's
neural language model in 2003, Mikolov's word2vec in 2013,
Vaswani's Transformer in 2017, and the rapid cascade of
scaled models that followed. None of this came out of
nowhere. All of it was, in retrospect, predictable from
the trajectory.

What you hold now is not an expert's knowledge. It is
something more specific and, for most purposes, more useful:
a working model of the mechanisms that produce both the
capability and the failure. You can reason forward from
the mechanism. When something goes wrong, you can ask why.
When a vendor makes a claim, you can ask what it would
take to verify it. When a colleague is dazzled, you can
point, gently, at the specific and definable ways the
system is also limited.

\section*{What changes and what does not}
\addcontentsline{toc}{section}{What changes and what does not}

The technology will change. New architectures will emerge.
Scaling laws will bend at points we cannot predict. Models
will become cheaper, larger, faster, stranger. Some of
what this guide has described will become obsolete within
years of its writing.

The mechanisms described here will not become obsolete
in the same way. The observation that next-token prediction
optimises for plausibility rather than truth is not a
property of any particular model; it is a property of the
objective. Until the objective changes, the observation
holds. The observation that a reward model trained on
ranked pairs learns a pattern-matcher rather than a judge
is not a property of RLHF or DPO specifically; it is a
property of any system that learns preferences from
pairwise comparisons. The observation that a function
approximated from text has no native way to connect form
to reference is Harnad's observation from 1990, and it
has not been resolved by the subsequent three decades of
scaling.

The honest posture, as we said in the preface, is this:
these systems are remarkable and limited in roughly equal
measure. Holding both halves simultaneously is the
beginning of useful thought about them.

\section*{Further reading}
\addcontentsline{toc}{section}{Further reading}

This guide was designed to be self-contained. For readers
who want to go further, the following are the most useful
next steps.

\medskip
\noindent\textbf{The technical companion.} \emph{Large
Language Models as Engineered Systems} is the technical
counterpart to this guide, covering the same thirteen
chapters with the formal mathematics included. Readers
who want to move from intuition to derivation, or who need
to evaluate vendor claims that rest on technical
specifications, should start there. The two books share
a chapter ordering and most of their historical narrative.

\medskip
\noindent\textbf{The original papers.} The field has been
generous about making its primary literature openly
accessible. Three papers are particularly valuable as
starting points.

Vaswani, A., Shazeer, N., Parmar, N., et al. (2017).
\emph{Attention Is All You Need}. The paper that introduced
the Transformer architecture. Available at arXiv:1706.03762.
Readable by a careful non-specialist; the intuition behind
the scaled dot-product attention mechanism is explained
well in Section~3.

Ouyang, L., Wu, J., Jiang, X., et al. (2022).
\emph{Training language models to follow instructions
with human feedback}. The InstructGPT paper, which laid
out the three-stage RLHF pipeline in detail. Available
at arXiv:2203.02155. The failure modes discussion in
the paper's limitations section is candid and worth reading.

Schaeffer, R., Miranda, B., and Koyejo, S. (2023).
\emph{Are Emergent Abilities of Large Language Models
a Mirage?} The paper that substantially shifted the
field's view of apparent emergent capabilities. Available
at arXiv:2304.15004. Short and accessible; the argument
about measurement artifacts in Section~2 is the core.

\medskip
\noindent\textbf{For the broader context.} Two books
serve as useful complements to the technical material.

Melanie Mitchell's \emph{Artificial Intelligence: A Guide
for Thinking Humans} (2019) covers the longer arc of the
field --- from the 1950s to the eve of the current wave ---
with a steady scepticism about strong claims in either
direction. It ages well because it is more interested
in the structural reasons for the field's failures than
in the specifics of any particular system.

Brian Christian's \emph{The Alignment Problem} (2020)
covers the alignment questions that Chapter~9 of this
guide treats. Christian's book is longer and more
detailed than our treatment, and it is particularly good
on the human side of the preference-data pipeline ---
the annotators, the rubrics, and the people who wrote them.

\medskip
\noindent\textbf{For staying current.} The field moves
faster than any book can follow. The most reliable
sources of current technical commentary, accessible to
a non-specialist, are the blogs maintained by the major
research laboratories (Anthropic, Google DeepMind,
OpenAI, Meta AI) and the summary newsletters that track
new papers. The important caveat is that these sources
are produced by parties with interests in the field's
reception. Read them for information; read them critically
for interpretation.

\section*{A final word}
\addcontentsline{toc}{section}{A final word}

You did not have to read this far. The fact that you did
suggests that you believe understanding these systems
matters --- that the right response to a technology that
has embedded itself in professional life at this speed
is not to treat it as magic, and not to dismiss it as
mere autocomplete, but to look at it clearly, with the
best vocabulary available.

That is the right response. The vocabulary is now yours.
What you do with it is the part that was always going to
be up to you.
